% GENERATED FILE. Edit canonical lessons, then run npm run generate:books.
% canonical-source-hash: fnv1a64:a3719e892163e835
% canonical-lessons: TA-C531-kutaippantu, TA-C531-irakuppantu, TA-C531-kaippantu, TA-C531-tennis, TA-C531-kapati

\chapter{Ball Games}
\label{ch:ta-ball-games}
\hlchaptermodality{531}

\begin{chapteropening}
\textbf{By the end of this chapter:} \emph{I can name ball games.}

Complete the last lesson of chapter 531: I can name ball games.
\end{chapteropening}

\section[\texorpdfstring{\ta{கூடைப்பந்து} (kūṭaippantu)}{கூடைப்பந்து (kūṭaippantu)}]{\ta{கூடைப்பந்து} (kūṭaippantu) — basketball}

\label{lesson:TA-C531-kutaippantu}



\begin{quote}
Before the new one: say the Tamil for carrom, then the Tamil for a pinwheel.
\end{quote}

\subsection*{You'll want to know: \ta{கூடைப்பந்து}}
\textbf{\ta{கூடைப்பந்து}} — \emph{kūṭaippantu} — \textquotedblleft{}basketball\textquotedblright{}.

\begin{quote}
\ta{எங்கள்} \ta{பள்ளியில்} \ta{கூடைப்பந்து} \ta{மைதானம்} \ta{இருக்கிறது}.
\end{quote}

\begin{quote}
\emph{eṅkaḷ paḷḷiyil kūṭaippantu maitāṉam irukkiṟatu.}
\end{quote}

\begin{quote}
Our school has a basketball court.
\end{quote}

Word for word \textquotedblleft{}basket-ball\textquotedblright{}: \ta{கூடை} (\emph{kūṭai}) + \ta{பந்து} (\emph{pantu}), with the \ta{ப்} doubled.

\begin{cousinweb}
\textbf{\ta{கூடைப்பந்து}} (\emph{kūṭaippantu}) is a native compound made to match English \emph{basketball}: \ta{கூடை} (\emph{kūṭai}), \textquotedblleft{}a basket\textquotedblright{} (DEDR 1884), and \ta{பந்து} (\emph{pantu}), \textquotedblleft{}a ball\textquotedblright{} (DEDR 3925).
\end{cousinweb}

\begin{grammarlens}[title={What you've built}]
One more word for describing what you lived and felt.
\end{grammarlens}

\subsection*{Your turn}
\begin{itemize}
\raggedright
  \item \emph{Say it:} \emph{kūṭaippantu}
  \item \emph{Say it:} \emph{kūṭaippantu}, once more
  \item \emph{Say it:} \emph{eṅkaḷ paḷḷiyil kūṭaippantu maitāṉam irukkiṟatu.} — Our school has a basketball court.
  \item \emph{Recall:} say the Tamil for carrom, then the Tamil for a pinwheel, then say \emph{kūṭaippantu} again
\end{itemize}

\begin{tcolorbox}[breakable,colback=teal!4,colframe=teal!35!black,title={Before you move on}]
What does \ta{கூடைப்பந்து} mean? (\textquotedblleft{}Basketball\textquotedblright{}.) Say it once more.
\end{tcolorbox}
\section[\texorpdfstring{\ta{இறகுப்பந்து} (iṟakuppantu)}{இறகுப்பந்து (iṟakuppantu)}]{\ta{இறகுப்பந்து} (iṟakuppantu) — badminton}

\label{lesson:TA-C531-irakuppantu}



\begin{quote}
Before the new one: say the Tamil for a pinwheel, then the Tamil for basketball.
\end{quote}

\subsection*{You'll want to know: \ta{இறகுப்பந்து}}
\textbf{\ta{இறகுப்பந்து}} — \emph{iṟakuppantu} — \textquotedblleft{}badminton\textquotedblright{}.

\begin{quote}
\ta{அக்காவும்} \ta{நானும்} \ta{மாலையில்} \ta{இறகுப்பந்து} \ta{விளையாடுகிறோம்}.
\end{quote}

\begin{quote}
\emph{akkāvum nāṉum mālaiyil iṟakuppantu viḷaiyāṭukiṟōm.}
\end{quote}

\begin{quote}
My elder sister and I play badminton in the evening.
\end{quote}

Word for word \textquotedblleft{}feather-ball\textquotedblright{}: \ta{இறகு} (\emph{iṟaku}) + \ta{பந்து}, after the shuttlecock.

\begin{cousinweb}
\textbf{\ta{இறகுப்பந்து}} (\emph{iṟakuppantu}) is a native compound: \ta{இறகு} (\emph{iṟaku}), \textquotedblleft{}a wing, a feather\textquotedblright{} (DEDR 2591), and \ta{பந்து} (\emph{pantu}), \textquotedblleft{}a ball\textquotedblright{} (DEDR 3925).
\end{cousinweb}

\begin{grammarlens}[title={What you've built}]
One more word for describing what you lived and felt.
\end{grammarlens}

\subsection*{Your turn}
\begin{itemize}
\raggedright
  \item \emph{Say it:} \emph{iṟakuppantu}
  \item \emph{Say it:} \emph{iṟakuppantu}, once more
  \item \emph{Say it:} \emph{akkāvum nāṉum mālaiyil iṟakuppantu viḷaiyāṭukiṟōm.} — My elder sister and I play badminton in the evening.
  \item \emph{Recall:} say the Tamil for a pinwheel, then the Tamil for basketball, then say \emph{iṟakuppantu} again
\end{itemize}

\begin{tcolorbox}[breakable,colback=teal!4,colframe=teal!35!black,title={Before you move on}]
What does \ta{இறகுப்பந்து} mean? (\textquotedblleft{}Badminton\textquotedblright{}.) Say it once more.
\end{tcolorbox}
\section[\texorpdfstring{\ta{கைப்பந்து} (kaippantu)}{கைப்பந்து (kaippantu)}]{\ta{கைப்பந்து} (kaippantu) — volleyball}

\label{lesson:TA-C531-kaippantu}



\begin{quote}
Before the new one: say the Tamil for basketball, then the Tamil for badminton.
\end{quote}

\subsection*{You'll want to know: \ta{கைப்பந்து}}
\textbf{\ta{கைப்பந்து}} — \emph{kaippantu} — \textquotedblleft{}volleyball\textquotedblright{}.

\begin{quote}
\ta{கடற்கரையில்} \ta{இளைஞர்கள்} \ta{கைப்பந்து} \ta{விளையாடினார்கள்}.
\end{quote}

\begin{quote}
\emph{kaṭaṟkaraiyil iḷaiñarkaḷ kaippantu viḷaiyāṭiṉārkaḷ.}
\end{quote}

\begin{quote}
Young people played volleyball on the beach.
\end{quote}

Word for word \textquotedblleft{}hand-ball\textquotedblright{}: \ta{கை} (\emph{kai}) + \ta{பந்து}.

\begin{cousinweb}
\textbf{\ta{கைப்பந்து}} (\emph{kaippantu}) is a native compound: \ta{கை} (\emph{kai}), \textquotedblleft{}a hand\textquotedblright{}, and \ta{பந்து} (\emph{pantu}), \textquotedblleft{}a ball\textquotedblright{} (DEDR 3925).
\end{cousinweb}

\begin{grammarlens}[title={What you've built}]
One more word for describing what you lived and felt.
\end{grammarlens}

\subsection*{Your turn}
\begin{itemize}
\raggedright
  \item \emph{Say it:} \emph{kaippantu}
  \item \emph{Say it:} \emph{kaippantu}, once more
  \item \emph{Say it:} \emph{kaṭaṟkaraiyil iḷaiñarkaḷ kaippantu viḷaiyāṭiṉārkaḷ.} — Young people played volleyball on the beach.
  \item \emph{Recall:} say the Tamil for basketball, then the Tamil for badminton, then say \emph{kaippantu} again
\end{itemize}

\begin{tcolorbox}[breakable,colback=teal!4,colframe=teal!35!black,title={Before you move on}]
What does \ta{கைப்பந்து} mean? (\textquotedblleft{}Volleyball\textquotedblright{}.) Say it once more.
\end{tcolorbox}
\section[\texorpdfstring{\ta{டென்னிஸ்} (ṭeṉṉis)}{டென்னிஸ் (ṭeṉṉis)}]{\ta{டென்னிஸ்} (ṭeṉṉis) — tennis}

\label{lesson:TA-C531-tennis}



\begin{quote}
Before the new one: say the Tamil for badminton, then the Tamil for volleyball.
\end{quote}

\subsection*{You'll want to know: \ta{டென்னிஸ்}}
\textbf{\ta{டென்னிஸ்}} — \emph{ṭeṉṉis} — \textquotedblleft{}tennis\textquotedblright{}.

\begin{quote}
\ta{அவள்} \ta{வாரம்} \ta{மூன்று} \ta{முறை} \ta{டென்னிஸ்} \ta{விளையாடுகிறாள்}.
\end{quote}

\begin{quote}
\emph{avaḷ vāram mūṉṟu muṟai ṭeṉṉis viḷaiyāṭukiṟāḷ.}
\end{quote}

\begin{quote}
She plays tennis three times a week.
\end{quote}

Written with the final \ta{ஸ்} of the English word.

\begin{cousinweb}
\textbf{\ta{டென்னிஸ்}} (\emph{ṭeṉṉis}) is a loan from English \emph{tennis}.
\end{cousinweb}

\begin{grammarlens}[title={What you've built}]
One more word for describing what you lived and felt.
\end{grammarlens}

\subsection*{Your turn}
\begin{itemize}
\raggedright
  \item \emph{Say it:} \emph{ṭeṉṉis}
  \item \emph{Say it:} \emph{ṭeṉṉis}, once more
  \item \emph{Say it:} \emph{avaḷ vāram mūṉṟu muṟai ṭeṉṉis viḷaiyāṭukiṟāḷ.} — She plays tennis three times a week.
  \item \emph{Recall:} say the Tamil for badminton, then the Tamil for volleyball, then say \emph{ṭeṉṉis} again
\end{itemize}

\begin{tcolorbox}[breakable,colback=teal!4,colframe=teal!35!black,title={Before you move on}]
What does \ta{டென்னிஸ்} mean? (\textquotedblleft{}Tennis\textquotedblright{}.) Say it once more.
\end{tcolorbox}
\section[\texorpdfstring{\ta{கபடி} (kapaṭi)}{கபடி (kapaṭi)}]{\ta{கபடி} (kapaṭi) — kabaddi}

\label{lesson:TA-C531-kapati}



\begin{quote}
Before the new one: say the Tamil for volleyball, then the Tamil for tennis.
\end{quote}

\subsection*{You'll want to know: \ta{கபடி}}
\textbf{\ta{கபடி}} — \emph{kapaṭi} — \textquotedblleft{}kabaddi\textquotedblright{}.

\begin{quote}
\ta{கிராமத்துத்} \ta{திருவிழாவில்} \ta{கபடிப்} \ta{போட்டி} \ta{நடந்தது}.
\end{quote}

\begin{quote}
\emph{kirāmattut tiruviḻāvil kapaṭip pōṭṭi naṭantatu.}
\end{quote}

\begin{quote}
There was a kabaddi match at the village festival.
\end{quote}

A team game: a raider crosses into the other side, touches players and must get back on one breath.

\begin{cousinweb}
The origin of \textbf{\ta{கபடி}} (\emph{kapaṭi}) is not settled: no dictionary we can check gives its source.
\end{cousinweb}

\begin{grammarlens}[title={What you've built}]
One more word for describing what you lived and felt.
\end{grammarlens}

\subsection*{Your turn}
\begin{itemize}
\raggedright
  \item \emph{Say it:} \emph{kapaṭi}
  \item \emph{Say it:} \emph{kapaṭi}, once more
  \item \emph{Say it:} \emph{kirāmattut tiruviḻāvil kapaṭip pōṭṭi naṭantatu.} — There was a kabaddi match at the village festival.
  \item \emph{Recall:} say the Tamil for volleyball, then the Tamil for tennis, then say \emph{kapaṭi} again
\end{itemize}

\begin{tcolorbox}[breakable,colback=teal!4,colframe=teal!35!black,title={Before you move on}]
What does \ta{கபடி} mean? (\textquotedblleft{}Kabaddi\textquotedblright{}.) Say it once more.
\end{tcolorbox}
